\documentclass{article}\usepackage[utf8]{inputenc}\usepackage[T1]{fontenc}\usepackage[french]{babel}\usepackage{geometry}\geometry{margin=2.5cm}\begin{document}\section*{Vérification de l'efficacité de la protection - MSAP}\subsection*{1. Diagnostic des habilitations et autorisations (Documents 1 et 2)}
Les tests effectués valident la pertinence des nouvelles restrictions mises en place :
\begin{itemize}
\item \textbf{Étanchéité des accès aux données} : Le compte "Enedis" accède normalement à ses fichiers clients , tandis que la tentative de connexion du compte "MSA" au même dossier est bloquée par un message d'erreur d'autorisation. Les permissions NTFS sont donc effectives.
\item \textbf{Contrôle de l'installation logicielle} : La tentative d'installation de PuTTY par l'utilisateur "Enedis" déclenche une demande de mot de passe administrateur. Cela confirme que les utilisateurs ne possèdent plus de privilèges élevés sur leurs sessions locales.
\end{itemize}\subsection*{2. Solution de renforcement du SI (Documents 3 et 4)}
Pour améliorer la surveillance et la réactivité face aux incidents, une solution complémentaire est nécessaire :
\begin{itemize}
\item \textbf{Centralisation des journaux (Logs)} : Bien que les événements de sécurité soient enregistrés localement sur chaque poste , il est souhaitable d'utiliser le dossier "Abonnements" de l'observateur d'événements du serveur d'administration.
\item \textbf{Objectif} : Transférer les logs des postes clients vers le serveur de M. Jivon pour permettre une analyse globale et détecter plus rapidement des tentatives d'intrusion ou des erreurs d'authentification massives.
\end{itemize}\subsection*{3. Efficacité de l'infrastructure physique et logique (Documents 5, 6 et 7)}
L'analyse des tests réseau montre que les objectifs de segmentation sont \textbf{partiellement} atteints :
\begin{itemize}
\item \textbf{Accès au serveur de fichiers} : L'objectif est rempli. Le \texttt{ping} et le \texttt{tracert} depuis un poste Enedis vers le serveur (192.168.100.1) sont réussis. L'ACL "ctsrv" autorise explicitement ce flux.
\item \textbf{Isolation entre partenaires} : L'objectif n'est \textbf{pas totalement rempli} selon les outils utilisés :
\begin{itemize}
\item Le \texttt{ping} vers un autre partenaire échoue avec le message "Destination host unreachable" , et le simulateur indique que le paquet est rejeté par la règle \texttt{deny ip any any}.
\item Cependant, la commande \texttt{tracert} vers 192.168.20.1 (MSA) semble aboutir dans certains tests.
\end{itemize}
\item \textbf{Anomalie de configuration} : Le document 5 montre que les ports \texttt{Fa0/1} et \texttt{Fa0/2} sont bien assignés aux VLAN 10 et 20. Toutefois, l'ACL est appliquée en \texttt{in} sur les interfaces sub-ethernet , ce qui bloque le trafic entrant vers le routeur mais peut nécessiter un affinage pour garantir une isolation totale et sans exception entre les sous-réseaux IP.
\end{itemize}\end{document}